\documentclass[11pt,a4paper]{article}
\usepackage[utf8]{inputenc}
\usepackage[T1]{fontenc}
\usepackage[french]{babel}
\usepackage[margin=2.5cm]{geometry}
\usepackage{enumitem}

\title{Ordre du jour de la réunion du \rule{3cm}{0.4pt}}
\author{Solène Delannoy-Pavy}
\date{}

\begin{document}
\maketitle

\section*{1. Modifications du plan depuis la dernière réunion}

\begin{itemize}
\item \textbf{Suppression du formalisme des MDP décomposables à action combinatoire.} Nous souhaitons mettre le
problème RTE au centre de l'article. Or ce formalisme s'écarte du problème RTE sur deux points. Le coût n'est pas
séparable, et les jours sans pénalité sont des variables globales.
\item \textbf{Lien avec les modèles graphiques probabilistes établi dès le début de la section~3.} Il permet de
montrer que toutes les approximations sont construites selon le même schéma.
\begin{enumerate}
\item On définit des groupes de variables (clusters) et des marginales sur ces groupes, qui sont les variables du
programme linéaire.
\item On approche le coût par une fonction linéaire de ces marginales, ce qui donne l'objectif du programme linéaire.
\item On remplace le polytope marginal (les marginales atteignables par les politiques du MDP) par le polytope local
(les marginales localement cohérentes), ce qui donne l'ensemble admissible du programme linéaire.
\end{enumerate}
\item \textbf{Résultats numériques uniquement sur le problème RTE.}
\end{itemize}

\section*{2. Questions sur la version actuelle}

\begin{enumerate}
\item \textbf{Rédaction.} Globalement, la rédaction actuelle me semble lourde. Comment l'améliorer ?
\item \textbf{Résultat sur les MDP décomposables.} Faut-il garder le résultat selon lequel l'approximation fluide peut
être arbitrairement mauvaise pour ces MDP ? Il s'éloigne un peu du problème RTE.
\item \textbf{Modèle simplifié.} Les résultats sur le modèle simplifié restent intéressants. Le modèle étant plus
simple, ses résultats sont plus faciles à interpréter. Faut-il les conserver ?
\item \textbf{Choix des reviewers.} Mettre en avant les modèles graphiques probabilistes risque d'attirer des
reviewers différents. Est-ce un problème ?
\item \textbf{Baseline.} Nous comparons nos formulations à une heuristique en deux temps.
\begin{itemize}
\item Hors ligne, on optimise les politiques de maintenance composant par composant.
\item En ligne, on choisit le planning de coût minimal qui réalise les maintenances données par ces politiques.
\end{itemize}
Cette heuristique ne s'étend pas facilement au cas avec quotas. Dans les tests actuels, on optimise les maintenances
par composant sans tenir compte des quotas, puis on attribue les quotas en ligne de manière gloutonne, pour couvrir
tous les jours de maintenance jusqu'à épuisement du quota. Cette baseline est-elle suffisante ?
\item \textbf{Light fluid LP.} Faut-il présenter la formulation avec un sous-système supplémentaire qui représente les
quotas (Light fluid LP) ? Elle est plus légère, mais moins performante que notre formulation principale, le Full
fluid LP. Sur une sous-station, elle est 7 fois plus rapide par décision pour un coût 10\,\% plus élevé.
\end{enumerate}

\end{document}
